../cictikz/src/cictikz/data/tex/cictikz_preamble.tex

% One-bit resistor-string DAC: three equal resistors from V_REF to ground,
% and a pair of pass devices that pick one of the two taps.
%
% b_0 selects the upper tap at 2/3 V_REF and b_0-bar the lower one at
% 1/3 V_REF, which is the case analysis printed beside the figure in
% l04_dac.md. The original marks the devices as plain NMOS switches.

\begin{circuitikz}[american, thick, transform shape, circuit ee IEC]

  ../cictikz/src/cictikz/data/tex/cictikz_lib.tex
  ../cictikz/src/cictikz/data/tex/rdac_lib.tex

  \def\xs{0}          % the resistor string
  \def\xo{3.2}        % the output rail

  % ---- The string, ground at the bottom and V_REF at the top ----
  % \vresistor is 1.5 tall; the extra half unit spaces the taps far enough
  % apart for a switch and its control label to sit between them.
  \draw (\xs,0.5) \vresistor{$R$} -- ++(0,0.5) coordinate (tb)
                  \vresistor{$R$} -- ++(0,0.5) coordinate (ta)
                  \vresistor{$R$} -- ++(0,0.5) coordinate (vtop);
  \draw (\xs,0.5) -- (\xs,0.2) \vground;

  \draw (vtop) -- ++(-1.0,0) coordinate (vrefport);
  \rdacPort{(vrefport)}{south}{$V_{REF}$}

  % ---- The two pass devices, tap on the left and output on the right ----
  \rdacSw{MA}{(ta)}{2.0}{$b_0$}
  \rdacSw{MB}{(tb)}{2.0}{$\overline{b_0}$}

  % ---- Both devices land on the output rail ----
  % Off the tap coordinates, not off numbers: \vresistor is \grid tall, so
  % the tap spacing is 2.1 rather than the 2.0 the layout reads like.
  \coordinate (ra) at ($(ta)+(2.0,0)$);
  \coordinate (rb) at ($(tb)+(2.0,0)$);
  \draw (ra) -- (\xo,0 |- ra) -- (\xo,0 |- rb) -- (rb);
  \coordinate (vo) at ($(\xo,0 |- ra)!0.5!(\xo,0 |- rb)$);
  \fill (vo) circle (0.07);
  \draw (vo) -- ++(0.8,0) coordinate (voutport);
  \rdacPort{(voutport)}{south west}{$V_{OUT}$}

\end{circuitikz}
\end{document}
